\subsubsection{Giải chi tiết Bài tập b (Phương pháp đơn hình)}

\textbf{1. Đưa bài toán về dạng chuẩn tắc}

Vì tất cả các ràng buộc chính đều mang dấu "$\le$", ta cộng thêm vào vế trái của mỗi ràng buộc một biến phụ không âm $x_4, x_5, x_6 \ge 0$ để tạo thành đẳng thức[cite: 163, 164, 2048, 2050, 2058, 2059, 2060, 2061].
Mô hình bài toán dạng chuẩn được viết lại như sau:

\begin{alignat*}{4}
    f(x) = \quad & 4x_1 & {}+{} & 2x_2 & {}+{} & 6x_3 & {}+{} 0x_4 + 0x_5 + 0x_6 \rightarrow \max \\
    \text{s.t.} \quad & x_1 & {}+{} & x_2 & {}+{} & 2x_3 & {}+{} x_4 = 8000 \\
    & x_1 & {}+{} & x_2 & {}+{} & 3x_3 & {}+{} x_5 = 16000 \\
    & 3x_1 & {}+{} & x_2 & {}+{} & 2x_3 & {}+{} x_6 = 12000 \\
    & \rlap{$x_j \ge 0, \quad j = \overline{1,6}$}
\end{alignat*}

Hệ vector điều kiện của các biến phụ $x_4, x_5, x_6$ tạo thành một ma trận đơn vị $I_3$, nên ta chọn $(x_4, x_5, x_6)$ làm cơ sở xuất phát.

\vspace{0.3cm}
\textbf{2. Quá trình giải bằng Bảng đơn hình}

\textbf{Bảng 1 (Bước lặp 0):}
\begin{table}[H]
    \centering
    \renewcommand{\arraystretch}{1.3}
    \begin{tabular}{|c|c|c|c|c|c|c|c|c|c|}
        \hline
        \multirow{2}{*}{\textbf{Cơ sở}} & \multirow{2}{*}{$\mathbf{C_b}$} & \multirow{2}{*}{\textbf{PA}} & $x_1$ & $x_2$ & $x_3$ & $x_4$ & $x_5$ & $x_6$ & \multirow{2}{*}{\textbf{Tỷ số $\lambda$}} \\
        \cline{4-9}
        & & & 4 & 2 & 6 & 0 & 0 & 0 & \\
        \hline
        $x_4$ & 0 & 8000 & 1 & 1 & \textbf{2} & 1 & 0 & 0 & $8000/2 = \mathbf{4000} \rightarrow$ \\
        \hline
        $x_5$ & 0 & 16000 & 1 & 1 & 3 & 0 & 1 & 0 & $16000/3 \approx 5333.3$ \\
        \hline
        $x_6$ & 0 & 12000 & 3 & 1 & 2 & 0 & 0 & 1 & $12000/2 = 6000$ \\
        \hline
        \multicolumn{2}{|c|}{$\Delta_j$} & \textbf{Z = 0} & -4 & -2 & \textbf{-6} $\uparrow$ & 0 & 0 & 0 & \\
        \hline
    \end{tabular}
\end{table}
\textit{Giải thích xoay bảng 1:}
\begin{itemize}[itemsep=0pt, parsep=0pt, topsep=0pt]
    \item Ta có $\Delta_3 = -6$ là số âm nhỏ nhất $\Rightarrow$ Chọn biến $x_3$ đưa \textbf{vào cơ sở}.
    \item Xét $\min\{8000/2, 16000/3, 12000/2\} = 4000$ tại dòng $x_4$ $\Rightarrow$ Chọn biến $x_4$ đưa \textbf{ra khỏi cơ sở}. Phần tử trụ là $\mathbf{2}$.
\end{itemize}

\vspace{0.3cm}
\textbf{Bảng 2 (Bước lặp 1):} Xoay bảng theo phần tử trụ $2$ bằng phép biến đổi Gauss-Jordan.
\begin{table}[H]
    \centering
    \renewcommand{\arraystretch}{1.3}
    \begin{tabular}{|c|c|c|c|c|c|c|c|c|c|}
        \hline
        \multirow{2}{*}{\textbf{Cơ sở}} & \multirow{2}{*}{$\mathbf{C_b}$} & \multirow{2}{*}{\textbf{PA}} & $x_1$ & $x_2$ & $x_3$ & $x_4$ & $x_5$ & $x_6$ & \multirow{2}{*}{\textbf{Tỷ số $\lambda$}} \\
        \cline{4-9}
        & & & 4 & 2 & 6 & 0 & 0 & 0 & \\
        \hline
        $x_3$ & 6 & 4000 & 1/2 & 1/2 & 1 & 1/2 & 0 & 0 & $4000 / (1/2) = 8000$ \\
        \hline
        $x_5$ & 0 & 4000 & -1/2 & -1/2 & 0 & -3/2 & 1 & 0 & - \\
        \hline
        $x_6$ & 0 & 4000 & \textbf{2} & 0 & 0 & -1 & 0 & 1 & $4000 / 2 = \mathbf{2000} \rightarrow$ \\
        \hline
        \multicolumn{2}{|c|}{$\Delta_j$} & \textbf{Z = 24000} & \textbf{-1} $\uparrow$ & 1 & 0 & 3 & 0 & 0 & \\
        \hline
    \end{tabular}
\end{table}
\textit{Giải thích xoay bảng 2:}
\begin{itemize}[itemsep=0pt, parsep=0pt, topsep=0pt]
    \item Bảng vẫn còn $\Delta_1 = -1 < 0$, nên phương án chưa tối ưu $\Rightarrow$ Chọn $x_1$ đưa \textbf{vào cơ sở}.
    \item Lập tỷ số $\lambda$ ở cột $x_1$ (chỉ tính dòng có hệ số $a_{i1} > 0$): $\min\{4000 / (1/2), 4000/2\} = 2000$ $\Rightarrow$ Chọn $x_6$ đưa \textbf{ra khỏi cơ sở}. Phần tử trụ là $\mathbf{2}$.
\end{itemize}

\vspace{0.3cm}
\textbf{Bảng 3 (Bước lặp 2):} Cập nhật bảng theo phần tử trụ $2$.
\begin{table}[H]
    \centering
    \renewcommand{\arraystretch}{1.3}
    \begin{tabular}{|c|c|c|c|c|c|c|c|c|}
        \hline
        \multirow{2}{*}{\textbf{Cơ sở}} & \multirow{2}{*}{$\mathbf{C_b}$} & \multirow{2}{*}{\textbf{PA}} & $x_1$ & $x_2$ & $x_3$ & $x_4$ & $x_5$ & $x_6$ \\
        \cline{4-9}
        & & & 4 & 2 & 6 & 0 & 0 & 0 \\
        \hline
        $x_3$ & 6 & 3000 & 0 & 1/2 & 1 & 3/4 & 0 & -1/4 \\
        \hline
        $x_5$ & 0 & 5000 & 0 & -1/2 & 0 & -7/4 & 1 & 1/4 \\
        \hline
        $x_1$ & 4 & 2000 & 1 & 0 & 0 & -1/2 & 0 & 1/2 \\
        \hline
        \multicolumn{2}{|c|}{$\Delta_j$} & \textbf{Z = 26000} & 0 & 1 & 0 & 5/2 & 0 & 1/2 \\
        \hline
    \end{tabular}
\end{table}

\textbf{3. Kết luận}

Tại Bảng 3, tất cả các hệ số ước lượng đều không âm ($\Delta_j \ge 0, \forall j$), do đó phương án hiện hành là phương án tối ưu.
\begin{itemize}
    \item Phương án tối ưu của bài toán là: $x^* = (2000, 0, 3000)$.
    \item Giá trị lớn nhất của hàm mục tiêu là: $f(x)_{\max} = 26000$.
\end{itemize}